\section{Descendant potentials and quantization}\label{conv:section}

We fix the ordinary theory and its operator normalization before introducing
the auxiliary equivariant target.  Quantized changes of frame are naturally
defined up to a scalar.  We therefore record both the exact Virasoro
constraints and the weaker, projective form that can be transported between
frames.  The scalar ambiguity will be removed at the end of the proof.

\subsection{Cohomology, curve classes, and descendants}

Let $Y$ be a smooth connected projective complex variety of dimension $d$.
Its state space is the full super vector space
$H_Y=H^*(Y;\C)$, with parity given by cohomological degree modulo two and
pairing
\[
 \eta_Y(a,b)=(a,b)_Y=\int_Y a\cup b.
\]
Every tensor product, permutation, contraction, and derivative below uses
the Koszul convention.  In particular, the coevaluation tensor is the
categorical inverse of this pairing, without an additional parity
involution.  Define even endomorphisms
\begin{equation}
 \mu_Y\big|_{H^{a,b}(Y)}=(a-d/2)\id,
 \qquad R_Y=c_1(TY)\cup.
 \label{conv:grading}
\end{equation}
Thus $\mu_Y$ uses the first Hodge index.  Poincar\'e duality and the Hodge
type of $c_1(TY)$ give
$\mu_Y^*=-\mu_Y$, $R_Y^*=R_Y$, and $[\mu_Y,R_Y]=R_Y$.

Choose a homogeneous Hodge basis $(\phi_a)$ with $\phi_0=1_Y$.
For $j\geq0$, let $t_j^a$ have the parity of $\phi_a$, and put
\[
 t_j=\sum_a t_j^a\phi_a,
 \qquad t(z)=\sum_{j\geq0}t_jz^j.
\]
These are formal supercommuting variables; the aggregate insertion $t_j$
is even.  The connected, ordinary, unreduced descendant potentials are
\begin{equation}
\begin{split}
 F_g^Y(t)
 &=\sum_{\substack{\beta\ \mathrm{effective}\\m\geq0}}
     \frac{Q^\beta}{m!}
     \int_{[\overline{\mathcal M}_{g,m}(Y,\beta)]^{\mathrm{vir}}}
       \prod_{i=1}^{m}
          \left(\sum_{j\geq0}\psi_i^j\ev_i^*t_j\right),\\
 Z_Y(t)&=\exp\left(\sum_{g\geq0}\hbar^{g-1}F_g^Y(t)\right).
\end{split}
\label{conv:potentials}
\end{equation}
Only stable maps contribute.  Here $\psi_i$ is the cotangent class at the
marked point of the stable map.  It is not an ancestor class.
The exponential includes disconnected domains and the empty domain.

The Novikov symbols retain the actual effective integral classes
$\beta\in H_2(Y;\Z)$, including $\beta=0$; multiplication is
$Q^{\beta_1}Q^{\beta_2}=Q^{\beta_1+\beta_2}$.
Fixing an ample integral class gives the Novikov completion: only finitely
many labels occur below each fixed ample degree.  All identities are read
coefficientwise in this completion, at finite order in the descendant
variables, and with the coefficientwise Laurent interpretation in
$\hbar$.  In particular, classes are not identified merely because their
divisor degrees agree.  The finiteness of effective homology labels of
bounded degree follows, for example, from the finite-type Chow spaces of
cycles of bounded degree in a projective embedding.

\subsection{The ordinary Virasoro operators}

On the loop space and its standard polarization use
\begin{equation}
\begin{split}
 \cH_Y&=H_Y((z^{-1})),
 \qquad
 \Omega_Y(f,g)=\Res_{z=0}(f(-z),g(z))_Y\,dz,\\
 \cH_Y&=H_Y[z]\oplus z^{-1}H_Y[[z^{-1}]].
\end{split}
\label{conv:polarization}
\end{equation}
The infinitesimal symplectic operators are
\begin{equation}
 \ell_{-1,Y}=z^{-1},\qquad
 \ell_{0,Y}=z\partial_z+\frac12+\mu_Y+\frac{R_Y}{z},\qquad
 \ell_{k,Y}=\ell_{0,Y}(z\ell_{0,Y})^k\quad(k\geq1).
 \label{conv:ordinary-modes}
\end{equation}
Their infinitesimal symplectic property follows from the adjoint identities
after \eqref{conv:grading}, including the sign of $z$ in the residue
pairing.

For an even infinitesimal symplectic operator $A$ on a paired loop space
with pairing $\eta$, write
\[
 \mathcal Q_A(f)=\tfrac12\Omega(f,Af).
\]
Use homogeneous Darboux coordinates for the displayed polarization, with
positive coordinates $q_j^a$ and dual negative coordinates $p_{j,a}$.
Normal ordering quantizes a $pp$ monomial as $\hbar$ times the
corresponding second derivative, a $pq$ monomial as the first-order
operator with multiplication before differentiation, and a $qq$ monomial
as multiplication by that monomial divided by $\hbar$.
Derivatives are left derivatives and all reorderings have their graded
signs.  We denote the resulting operator, after the dilaton translation
\begin{equation}
                  q(z)=t(z)-z1_Y,
                  \label{conv:dilaton}
\end{equation}
by $\op_\eta(A)$.  The negative coordinates $p_{j,a}$ in this paragraph
are unrelated to the tautological divisor $p$ introduced below.

With this sign convention the string operator is
\begin{equation}
 L_{-1}^Y=\op_{\eta_Y}(\ell_{-1,Y})
   =-\frac{\partial}{\partial t_0^0}
     +\sum_{j\geq0,a}t_{j+1}^a\frac{\partial}{\partial t_j^a}
     +\frac{(t_0,t_0)_Y}{2\hbar}.
 \label{conv:string}
\end{equation}
For all other modes set
\begin{equation}
\begin{split}
 L_k^Y&=\op_{\eta_Y}(\ell_{k,Y})\quad(k\ne0),\\
 L_0^Y&=\op_{\eta_Y}(\ell_{0,Y})+C_Y,
 \qquad
 C_Y=\frac{\chi(Y)}{16}-\frac14\str(\mu_Y^2).
\end{split}
\label{conv:normalized-modes}
\end{equation}
The supertrace uses even minus odd cohomological parity.  There are no
scalar corrections in positive modes.  We write $\V(Y)$ for the full set
of identities
\begin{equation}
                    L_k^Y Z_Y=0\qquad(k\geq-1).
                    \label{conv:virasoro}
\end{equation}
These equations use every genus, every effective integral class, and all
insertions in $H_Y$, including odd and primitive classes.

\subsection{Projective constraints and changes of frame}

\begin{definition}\label{conv:projective}
An even infinitesimal symplectic operator $A$ is \emph{satisfied
projectively} by a potential $Z$ if
\[
                       Z^{-1}\op_\eta(A)Z
\]
is independent of the descendant or ancestor variables.  In this
expression $\op_\eta(A)$ acts on $Z$.  Such a scalar may depend on the
Novikov variables, equivariant parameters, and $\hbar$.
\end{definition}

This formulation discards exactly the scalar ambiguity of quadratic
quantization.  It is useful for comparing different paired state spaces.
For a symplectic map $M$ between them, denote its quantized action, when
defined in the completions used here, by $U_M$.  Our group-action
convention is the one in the following covariance formula.  With the
Hamiltonian sign above, its infinitesimal form uses $-\op_\eta(A)$ for
the action of $\exp(A)$.

\begin{lemma}[Projective covariance]\label{conv:covariance}
Suppose $M$ and $M^{-1}$ admit quantized actions, and $A$ and $MAM^{-1}$
admit coefficientwise quadratic quantizations.  If $\eta$ and $\eta'$
are the source and target pairings, then
\begin{equation}
 U_M\op_\eta(A)U_M^{-1}
        =\op_{\eta'}(MAM^{-1})+\text{a scalar}.
 \label{conv:covariance-equation}
\end{equation}
Consequently $A$ is satisfied projectively by $Z$ if and only if
$MAM^{-1}$ is satisfied projectively by $U_M Z$.
Multiplying either potential by an invertible scalar does not change this
condition.
\end{lemma}

\begin{proof}
The commutator of two normally ordered quadratic operators is the
quantization of the corresponding quadratic Hamiltonian bracket, together
with the scalar arising from the double contractions.  In super
coordinates that scalar is the corresponding supertrace.  Exponentiating
this identity gives \eqref{conv:covariance-equation}; equivalently one can
use the projective quantization formalism of \cite{GiventalQuant}.
The same argument applies between paired spaces after a linear change of
Darboux coordinates.  The last assertions follow because these operators
do not differentiate coefficient parameters or $\hbar$.
\end{proof}

\subsection{Coefficientwise finiteness near zero curve variables}

Some later mode matrices have infinitely many positive powers of $z$,
and can also contain finite powers of $z\partial_z$.  We interpret their
residue identities entrywise in the polarized mode coordinates.  The
arrays need not act on every uncompleted loop-space input; only the
coefficientwise actions specified here are used.  The
following elementary rule specifies the formal setting for their local
quantizations; the continuation argument will establish its separate
finiteness assertion at generic fiber parameter.

\begin{lemma}\label{conv:finite-support}
Suppose a potential has finite descendant-index support at fixed curve
coefficient, $\hbar$ power, and variable order, and has $\hbar$ powers
bounded below at fixed curve coefficient and variable order.  Let $A$ be
an even infinitesimal symplectic mode operator, linear over the coefficient
parameters and independent of $\hbar$.  Assume its powers of $z$ have a
finite lower bound and its order in $z\partial_z$ is finite at each
filtration coefficient.  Then $\op_\eta(A)$ acts coefficientwise on this
potential.  The same holds with additional formal filtration variables.
Quantized actions of transformations equal to the identity modulo
positive filtration, whose logarithms satisfy these mode bounds, are
defined by their formal exponential series in this setting.
\end{lemma}

\begin{proof}
For the $pp$ part, only finitely many differentiated indices have nonzero
coefficients in the potential.  For the $pq$ part, a fixed differentiated
index bounds the possible multiplied index because the mode powers have
a lower bound.  The $qq$ part has only finitely many pairs of indices at
each filtration coefficient.  Euler differentiation in $z$ multiplies
mode coefficients by polynomials in their indices and does not change
these bounds.  At fixed positive filtration degree only finitely many
terms of an exponential occur.  These observations also make the
cross-polarization contractions in the covariance formula finite
coefficientwise.
\end{proof}

Ordinary descendant potentials have this support property: for fixed
genus, class, and number of marks, sufficiently high powers of the
cotangent classes vanish on the finite-dimensional stable-map space.
It also holds for the twisted descendants used below, where the torus
acts trivially on the target and the cotangent classes remain ordinary
classes.  Ancestor potentials have the stronger bound supplied by the
dimension of the moduli space of stable curves.
Stability ensures the stated Laurent property after exponentiation:
degree-zero genus-zero components consume marked points, whereas
positive-degree components consume positive Novikov degree.

The standard splitting and cotangent-class identities used in the paper
are tensor identities with the diagonal given by the categorical
coevaluation.  Their marked-point proofs therefore apply to full
cohomology with precisely these signs.  Characteristic-class
multiplications are even; quantum Riemann--Roch and the two-mark shift
identities likewise use arbitrary evaluation classes and this diagonal.
Throughout, a splitting adds the actual integral curve classes.  These
conventions permit the standard formulas to be used without discarding
odd insertions or merging Novikov labels.
